\section{Discussion}
\label{sec:discussion}
The tests show how the configured runtime handled the cases exercised here. Mapping probes exercised configured signal bindings, and the process-window ablation showed that the tested mock-driver route rejected the six selected out-of-window requests when the interlock was enabled. In the E6 comparison, all nine valid pairs had lower IAE under the deterministic recipe policy than under no-write. This is an exploratory result from a small simulated sample, not a general control-effect estimate. The earlier 15/15 versus 0/15 endpoint counts used different horizons and endpoint rules and remain historical descriptive evidence in Supplementary File~S7.

The three E7 episodes ended canceled and below the target band. Their local recorder reported two writes per episode, but the archive lacks a server-side journal for independent reconciliation. The task-mediated and direct routes share dispatch in the design, but route equivalence and cross-harness robustness were not tested. The single twin run also lacks raw samples and a frozen runtime snapshot. More broadly, the guard ablation had no readback-mismatch stimulus, approvals were scripted, and no PID baseline, operator study, or physical-plant evaluation was performed. A larger preregistered process comparison with complete source and event records is needed to estimate control effects; Agent execution and approval quality require separate task-level and human-operator studies. Simulator and mock-driver checks complement, but do not replace, PLC, safety instrumented system (SIS), and equipment protections.

\section{Conclusion}
\label{sec:conclusion}
The runtime binds production context and direction-specific mappings to shared dispatch, then applies route-specific checks before a write reaches a registered driver. It also keeps command outcomes, optional readbacks, and later process observations as separate records. In the tested simulator configurations, mapping probes exercised selected bindings, the guard ablation rejected the targeted window violations when the interlock was enabled, and nine valid E6 pairs showed lower IAE under the deterministic scripted policy than under no-write. The E6 result is exploratory; the three E7 episodes were canceled and did not reach the target band. These results do not establish general control efficacy, reliable autonomous task execution, operator decision quality, or physical-plant safety. Larger controlled process comparisons and complete Agent and operator studies remain open.
